\documentclass{article}
\usepackage[UTF8]{ctex}
\usepackage{amsmath,amssymb,amsfonts}
\usepackage{enumitem}
\usepackage{geometry}
\geometry{a4paper, margin=1in}
% === 错题本排版优化（不参与题目标识） ===
\setlength{\parindent}{0pt}
\setlength{\parskip}{0.35em plus 0.10em minus 0.08em}
\setlength{\emergencystretch}{3em}
\setlist{itemsep=0.25em, topsep=0.40em, parsep=0pt}
\allowdisplaybreaks[2]
% === 错题本排版优化结束 ===
\begin{document}
\section*{南昌大学2024年高等代数真题}
\section*{111．南昌大学2024年高等代数试卷}

\subsection{题目1}
1．已知多项式 \(f(x)=1+x+\cdots+x^{n-1}\) ，证明：\(f(x) \mid\left[f(x)+x^{n}\right]^{2}-x^{n}\) ．

\subsection{题目2}
2．计算 \(n\) 阶行列式

\[
D_{n}=\left|\begin{array}{ccccc}
x_{1}^{2}-2 & x_{1} x_{2} & x_{1} x_{3} & \cdots & x_{1} x_{n} \\
x_{2} x_{1} & x_{2}^{2}-2 & x_{2} x_{3} & \cdots & x_{2} x_{n} \\
x_{3} x_{1} & x_{3} x_{2} & x_{3}^{2}-2 & \cdots & x_{3} x_{n} \\
\vdots & \vdots & \vdots & & \vdots \\
x_{n} x_{1} & x_{n} x_{2} & x_{n} x_{3} & \cdots & x_{n}^{2}-2
\end{array}\right|
\]

\subsection{题目3}
3．设 \(\alpha_{i}=\left(a_{i 1}, a_{i 2}, \cdots, a_{i n}\right)(i=1,2, \cdots, s), \beta=\left(b_{1}, b_{2}, \cdots, b_{n}\right)\) 。证明：若线性方程组

\[
\left\{\begin{array}{c}
a_{11} x_{1}+a_{12} x_{2}+\cdots+a_{1 n} x_{n}=0 \\
a_{21} x_{1}+a_{22} x_{2}+\cdots+a_{2 n} x_{n}=0 \\
\cdots \cdots \\
a_{n 1} x_{1}+a_{n 2} x_{2}+\cdots+a_{n n} x_{n}=0
\end{array}\right.
\]

的解全是 \(b_{1} x_{1}+b_{2} x_{2}+\cdots+b_{n} x_{n}=0\) 的解，那么 \(\beta\) 可经 \(\alpha_{1}, \alpha_{2}, \cdots, \alpha_{s}\) 线性表出。

\subsection{题目4}
4．设矩阵 \(A=\left(a_{i j}\right)_{n \times n}\) 正定．证明：

\begin{itemize}
  \item[（1）] A 的主对角线上的元素全大于 0 。
  \item[（2）] \(\left|a_{i j}\right|<\frac{1}{2}\left(a_{i i}+a_{j j}\right), i \neq j ; i, j=1,2, \cdots, n\) ．
  \item[（3）] \(A\) 中各元素绝对值最大者一定在主对角线上．
\end{itemize}

\subsection{题目5}
5．证明：\(\left(A^{*}\right)^{*}=|A|^{n-2} A\) ，其中 \(A\) 为 \(n \times n\) 矩阵 \((n \geq 2), A^{*}\) 为 \(A\) 的伴随矩阵。

\subsection{题目6}
6．设二次型 \(f\left(x_{1}, x_{2}, x_{3}\right)=2 x_{1}^{2}+x_{2}^{2}+x_{3}^{2}+2 x_{1} x_{2}+t x_{2} x_{3}\) ．

\begin{itemize}
  \item[（1）] 求 \(t\) 使得二次型 \(f\left(x_{1}, x_{2}, x_{3}\right)\) 正定。
  \item[（2）] \(t=1\) 时，求二次型 \(f\left(x_{1}, x_{2}, x_{3}\right)\) 对应矩阵 \(A\) 的最小多项式 \(m_{A}(\lambda)\) ．
\end{itemize}

\subsection{题目7}
7．设 \(A\) 为 \(n\) 阶方阵，证明：

\begin{itemize}
  \item[（1）] \(A^{n}=O\) 的充要条件是 \(A\) 的特征值全为零。
  \item[（2）] 若 \(A^{n}=O\) ，则 \(|A+E|=1\) 。
\end{itemize}

\subsection{题目8}
8 ．令 \(A=\left(\begin{array}{cc}\frac{3}{2} & -\frac{1}{2} \\ \frac{1}{2} & \frac{1}{2}\end{array}\right)\) ，求 \(A^{100}\) ．

\subsection{题目9}
9．设 \(\sigma\) 是数域 \(P\) 上线性空间 \(V\) 的线性变换，\(f(\lambda), m(\lambda)\) 分别为 \(\sigma\) 的特征多项式与最小多项式。且

\[
f(\lambda)=(\lambda+1)^{3}(\lambda-2)^{2}(\lambda+3), m(\lambda)=(\lambda+1)^{2}(\lambda-2)(\lambda+3) .
\]

\begin{itemize}
  \item[（1）] 写出 \(\sigma\) 的所有不变因子．
  \item[（2）] 写出 \(\sigma\) 的若尔当标准形．
\end{itemize}

\subsection{题目10}
10．设 \(A\) 为 \(n\) 阶非零实矩阵 \((n \geq 3)\) ，满足 \(A^{\mathrm{T}}=A^{*}\) ，其中 \(A^{\mathrm{T}}\) 为 \(A\) 的转置，\(A^{*}\) 为 \(A\) 的伴随矩阵。证明：

\begin{itemize}
  \item[（1）] \(|A|>0\) ．
  \item[（2）] \(A\) 为正交矩阵．
\end{itemize}
\end{document}
